\sscp{Reliability of the sources}{02-04-03}
Comparative linguists’ handling of the data can be only as good
as the sources of their data,
some of which are already secondary sources,
and some of the primary sources may be mislabelled,
under-differentiated, actually represent another language or
variety, be poorly transcribed, or misglossed.

For example, \citet{BlustTrussel2020} list a large quantity of
data under the label of “Buruese”.\footnote{
		The label “Buruese”
		is not from the \ili{Buru} langauge, it is not Indonesian, and not even good \ili{English}.
		Old Dutch sources (24 publications by Hendriks,
		Jansen, Jellesma, Schut, van der Miesen, Wilken,
		Willer) refer to the island,
		people and language as “Boeroe”
		and “Boeroetaal”, with one instance of “Burusch”.
		Devin also used “\ili{Buru}” for his unpublished dictionary.
		“\ili{Buru}” suffices to identify the language,
		and is much preferred.}
 Much of that data is from an unpublished dictionary by Ch.
Devin who did not distinguish different varieties of \ili{Buru} or
the \ili{Lisela} language, but tended to lump everything together.
Devin’s main informant, Bapak Temi,
was from Banu Lalet,
near the headwaters where the \ili{Lisela},
\ili{Rana}, Wae Sama, and \ili{Masarete} varieties converge.
However, several forms cited as “Buruese” are identifiable by
form or meaning as actually being from the related \ili{Lisela} [ISO:
lcl] language -- which has some patterned differences with the
\ili{Buru} language [ISO: mhs] in both the phonology (e.g.
different handling of PMP *u,
different handling of PMP *p,
sometimes a different vowel in antepenultimate syllable,
sometimes different CC onsets; \citet[37--38]{Grimes1991c}
and the lexicon (lexical similarity of \ili{Lisela} with several \ili{Buru}
dialects averages only 60--63{\%}).
Again, this reflects on the source,
but ultimately muddies the historical-comparative work.

From 1895--1939 several versions of an approximately 1,000-item
word list known as the “Holle lists” were prepared
and sent throughout what is now the Indonesian archipelago.
\citet[17]{Stokhof1980a} notes, “It was planned that the lists
would be completed by Dutch civil servants,
officers, missionaries and ‘intelligente Inlanders’ such as village
heads, merchants and teachers, in practice all people with a reasonable
knowledge of Dutch.” A number of scholars have noted the unreliable
labelling of language varieties or transcriptions or glosses
in some Holle lists.
A common problem is that those who filled out the lists were
often from entirely different regions,
different sound systems, and different semantic ranges,
untrained in linguistics, showing goodwill in trying to help
collect the desired data.
We are grateful for that,
but some word lists and whole sections of the data can be notoriously unreliable.
\citet[10--11]{Grimes1991c} critiques several Holle lists associated
with \ili{Buru} and \ili{Kayeli}.
For example, the Holle list labelled “\ili{Hukumina}” \citep[63ff]{Stokhof1982b}
is not just the recently extinct \ili{Hukumina} language from western
\ili{Buru}, but also contains a mixture of the \ili{Lisela}
and \ili{Kayeli} languages from a village named \ili{Hukumina} near \ili{Kayeli}
village in eastern \ili{Buru} \citep{Grimes2010b}.

